\section*{Introducción}
Este texto tratará el análisis y diseño de circuitos con amplificadores lineales integrados. El capítulo I y II están dedicados al amplificador operacional 
clásico, modernamente identificado por los fabricantes como realimentado por 
tensión, (Voltage Feedback Amplifier), para diferenciarlo de otros tipos. \\
Únicamente se analizaran las aplicaciones de las configuraciones realimentadas, 
utilizando  en el análisis las técnicas y conclusiones de este tipo de sistemas. \\
En el  capítulo. 
I, se tratará el amplificador ideal a los efectos de familiarizase  
con técnicas operativas clásicas y definir magnitudes que serán referencias. 
Para el  estudio del capítulo II, que se centrará en el amplificador operacional 
real, con función de transferencia de un solo polo (Internamente compensado), 
donde se analizarán las desviaciones producidas por sus características reales, 
profundizando aquellas que se refieren a la respuesta en frecuencia del conjunto que son generalmente el punto crítico en la síntesis de amplificadores de 
banda ancha. En el capítulo III se estudiará otro tipo de amplificador de muy 
alta ganancia y velocidad, denominado  realimentado por corriente, (CFA) 
mientras que el IV, lo hará con los amplificadores de transconductancia, 
(O.T.A), en sus dos variantes (tecnología VFA y CFA). El capítulo V y parte del 
VI tratan temas generales respecto del modelado de la respuesta en frecuencia 
de redes activas  (Realimentadas o Directas), cuyas conclusiones serán de 
aplicación en la síntesis de amplificadores en lo que resta del texto.
\newpage